\chapter{Conclusion}
\label{chap:ch10}
\hspace{\parindent}
This dissertation addressed the problem of improving centralized multi-robot path planning in a
TEA*-based coordination framework. The work was motivated by the sensitivity of sequential
prioritized planning to robot ordering, where different priority permutations can lead to substantially
different outcomes in terms of feasibility, solution cost, and planning time. Since exhaustive
exploration of all possible robot orderings is infeasible for large fleets, the dissertation investigated
heuristic and parallel strategies for exploring more relevant priority orderings within a practical
planning framework.

The first part of the work established the theoretical and technical foundations required for this
study, including MAPF formulations, graph-based search, clustering, regression, and parallel
computing concepts. The existing CRIIS fleet coordination system was then analyzed as the baseline
architecture, with particular focus on its centralized planning pipeline and on the role of TEA* as the
underlying path planning method. This analysis showed that TEA* provides a practical mechanism
for computing time-aware, conflict-free plans, but that its sequential structure makes the quality of
the final solution strongly dependent on the selected robot priority ordering.

To address this limitation, several prioritization heuristics were designed and implemented. These
heuristics were organized into different families, namely \textit{Roadmap Distance},
\textit{Topological}, \textit{Surroundings}, \textit{Conflicts}, and \textit{Replanning}, according to the
type of information used to compute robot priority scores. Their behavior was evaluated using an
isolated Path Planner testing framework, which allowed the effect of each priority ordering to be
analyzed independently from the remaining coordination stack. The experimental results showed
that no single heuristic was consistently dominant across all scenarios. Instead, each heuristic family
was more effective under different spatial, topological, or conflict-related conditions, confirming the
importance of considering multiple priority orderings rather than relying on a fixed ordering rule.

Building on this observation, a parallel planning framework was developed to evaluate multiple
priority orderings concurrently. In this framework, worker threads execute independent TEA*
instances, each using a different heuristic-generated ordering, while shared mechanisms manage
heuristic dispatching, timeout handling, candidate plan selection, and redundant heuristic filtering.
This design preserves the underlying TEA* planner while extending the exploration of the priority
ordering space, allowing the system to search for better solutions without replacing the existing
centralized planning architecture.

The complete CRIIS coordination stack was also analyzed during integrated execution, where
temporal consistency issues were observed in the processing of planning requests and robot state
updates. These issues could lead to redundant replanning operations and planning decisions based on
incomplete temporal context. To mitigate this behavior, batching mechanisms were introduced for
both planning requests and robot state updates, reducing unnecessary replanning activity and
improving the consistency of the coordination pipeline during execution.

Overall, the results obtained in this dissertation show that the sensitivity of sequential TEA*-based
planning to robot ordering was mitigated in the evaluated scenarios without abandoning the existing
planner architecture. The proposed heuristics produced diverse and useful robot sequences, while
the parallel planner enabled several of these alternatives to be evaluated within the same planning
cycle. In addition, the batching mechanisms reduced redundant replanning activity and improved
temporal consistency during the observed integrated executions. Together, these contributions
provide a practical extension of the baseline centralized fleet coordination system, with observed
gains in valid-solution reliability, ordering flexibility, and suitability for the evaluated
large-scale planning scenarios.